\documentclass[border=12pt]{standalone}
\usepackage[no-math]{fontspec}
\usepackage[brazil]{babel}
\usepackage{amsmath,amssymb}
\IfFileExists{C:/Windows/Fonts/Lato_Regular.ttf}{%
  \setmainfont{Lato}[
    Path = C:/Windows/Fonts/,
    UprightFont = Lato_Regular.ttf,
    BoldFont = Lato_Bold.ttf,
    ItalicFont = Lato_Italic.ttf,
    BoldItalicFont = Lato_BoldItalic.ttf
  ]
  \setsansfont{Lato}[
    Path = C:/Windows/Fonts/,
    UprightFont = Lato_Regular.ttf,
    BoldFont = Lato_Bold.ttf,
    ItalicFont = Lato_Italic.ttf,
    BoldItalicFont = Lato_BoldItalic.ttf
  ]
}{%
  \RequirePackage[default]{lato}
}
\renewcommand{\familydefault}{\sfdefault}

\usepackage{xcolor}
\usepackage{tikz}
\usepackage{pgfplots}
\pgfplotsset{compat=1.18}
\usetikzlibrary{arrows.meta,positioning,shapes.geometric,fit,backgrounds,calc,decorations.pathreplacing,patterns}

% Paleta institucional UEFS / Extensão Rural
\definecolor{erNavy}{HTML}{1F4268}
\definecolor{erNavyDark}{HTML}{152E4A}
\definecolor{erGold}{HTML}{B38A3D}
\definecolor{erGoldDark}{HTML}{8C6927}
\definecolor{erTeal}{HTML}{0E7C7B}
\definecolor{erForest}{HTML}{2D6A4F}
\definecolor{erAlert}{HTML}{9E2A2B}
\definecolor{erInk}{HTML}{26313D}
\definecolor{erMuted}{HTML}{5B6573}

\definecolor{bgNavy}{HTML}{EDF2F7}
\definecolor{bgGold}{HTML}{F8F3E9}
\definecolor{bgTeal}{HTML}{EBF6F5}
\definecolor{bgForest}{HTML}{EEF5F0}
\definecolor{bgAlert}{HTML}{FDF2F2}

\tikzset{
  >=Stealth,
  erBox/.style={
    draw=#1, fill=white, rounded corners=5pt, line width=1.15pt,
    inner sep=7pt, align=center, font=\footnotesize\sffamily, text=erInk
  },
  erTint/.style={
    draw=#1, fill=#1!8!white, rounded corners=5pt, line width=1.15pt,
    inner sep=7pt, align=center, font=\footnotesize\sffamily, text=erInk
  },
  erBadge/.style={
    fill=#1, text=white, rounded corners=9pt, inner xsep=8pt, inner ysep=3pt,
    font=\scriptsize\bfseries\sffamily
  },
  erTitle/.style={font=\normalsize\bfseries\sffamily, text=erNavy, align=center},
  erSub/.style={font=\footnotesize\sffamily, text=erMuted, align=center},
  erArrow/.style={->, >=Stealth, line width=1.35pt},
  erDArrow/.style={<->, >=Stealth, line width=1.35pt}
}


\begin{document}
\begin{tikzpicture}

  % Título principal
  \node[erTitle, font=\large\bfseries\sffamily] (title) at (0, 5.55)
    {Comunicação Dialógica (Freire) e Ferramentas do Diagnóstico Rural Participativo (DRP)};

  % ================= CONTRASTE SUPERIOR: FREIRE =================
  \node[erTint=erAlert, minimum width=7.5cm, minimum height=1.95cm, text width=6.9cm, align=center]
    (invasao) at (-4.45, 3.8) {
    \textbf{\small Extensão como Invasão Cultural (Bancária)}\\[5pt]
    \scriptsize\color{erInk}
    Transmissão vertical e antidialógica de receitas técnicas.\\
    Ignora o contexto sociocultural e silencia o saber camponês.
  };

  \node[erTint=erForest, minimum width=7.5cm, minimum height=1.95cm, text width=6.9cm, align=center]
    (dialogo) at (4.45, 3.8) {
    \textbf{\small Comunicação Horizontal e\\Problematizadora (Freire, 1977)}\\[4pt]
    \scriptsize\color{erInk}
    Co-investigação crítica da realidade agrária.\\
    Educador e agricultor aprendem juntos em reciprocidade.
  };

  % Seta de transição entre os modelos
  \draw[erArrow, color=erForest, line width=1.6pt] (invasao.east) -- (dialogo.west)
    node[midway, above=4pt, font=\scriptsize\bfseries\sffamily, text=erForest] {superação};

  % ================= DIVISOR DE SEÇÃO: 4 FERRAMENTAS DRP =================
  \node[erBadge=erNavy, rounded corners=5pt, inner ysep=4pt] (drpHeader) at (0, 2.05)
    {Principais Ferramentas do Diagnóstico Rural Participativo (DRP)};

  % Barra distribuidora para os 4 cards
  \draw[draw=erNavy!50, line width=1.1pt] (drpHeader.south) -- (0, 1.48);
  \draw[draw=erNavy!50, line width=1.1pt] (-6.3, 1.48) -- (6.3, 1.48);

  % ================= 4 CARDS DRP EM LINHA (SEM TIKZ ANINHADO!) =================
  % Card 1: Mapa Falado
  \node[erBox=erNavy, minimum width=3.9cm, minimum height=3.45cm] (card1) at (-6.3, -0.65) {};
  \draw[erArrow, color=erNavy!75, line width=1.1pt] (-6.3, 1.48) -- (card1.north);
  \node[erBadge=erNavy, rounded corners=5pt, inner ysep=4pt]
    at ([yshift=-18pt]card1.north) {1. Mapa Falado};
  \node[text width=3.45cm, align=center, font=\scriptsize\sffamily, text=erInk]
    at ([yshift=-14pt]card1.center) {
    \textbf{Espaço e Recursos}\\[6pt]
    \color{erMuted}Desenho coletivo da comunidade: limites, águas, solos, lavouras e infraestrutura local.
  };

  % Card 2: Travessia / Caminhada
  \node[erBox=erTeal, minimum width=3.9cm, minimum height=3.45cm] (card2) at (-2.1, -0.65) {};
  \draw[erArrow, color=erTeal!80, line width=1.1pt] (-2.1, 1.48) -- (card2.north);
  \node[erBadge=erTeal, rounded corners=5pt, inner ysep=4pt]
    at ([yshift=-18pt]card2.north) {2. Travessia / Caminhada};
  \node[text width=3.45cm, align=center, font=\scriptsize\sffamily, text=erInk]
    at ([yshift=-14pt]card2.center) {
    \textbf{Observação \textit{in loco}}\\[6pt]
    \color{erMuted}Percurso transversal no agroecossistema dialogando sobre relevo, manejo e problemas.
  };

  % Card 3: Calendário Sazonal
  \node[erBox=erForest, minimum width=3.9cm, minimum height=3.45cm] (card3) at (2.1, -0.65) {};
  \draw[erArrow, color=erForest!80, line width=1.1pt] (2.1, 1.48) -- (card3.north);
  \node[erBadge=erForest, rounded corners=5pt, inner ysep=4pt]
    at ([yshift=-18pt]card3.north) {3. Calendário Sazonal};
  \node[text width=3.45cm, align=center, font=\scriptsize\sffamily, text=erInk]
    at ([yshift=-14pt]card3.center) {
    \textbf{Chuvas, Plantio e Renda}\\[6pt]
    \color{erMuted}Mapeamento anual de precipitação, safras, mão de obra, pragas e fluxo financeiro.
  };

  % Card 4: Diagrama de Venn
  \node[erBox=erGold, minimum width=3.9cm, minimum height=3.45cm] (card4) at (6.3, -0.65) {};
  \draw[erArrow, color=erGoldDark!80, line width=1.1pt] (6.3, 1.48) -- (card4.north);
  \node[erBadge=erGoldDark, rounded corners=5pt, inner ysep=4pt]
    at ([yshift=-18pt]card4.north) {4. Diagrama de Venn};
  \node[text width=3.45cm, align=center, font=\scriptsize\sffamily, text=erInk]
    at ([yshift=-14pt]card4.center) {
    \textbf{Atores e Instituições}\\[6pt]
    \color{erMuted}Círculos representando importância, vínculo e distância das organizações no território.
  };

\end{tikzpicture}
\end{document}
